
\begin{proof}
By \Cref{thm:topological-family-closure}, the family section is the rigid four-punctured sphere
\[
X_f^\circ\cong\mathbb P^1\setminus\mu_4,
\]
and it carries the full $D_4$ symmetry inherited from the seam winding data. The determinant line
is fixed to be trivial, and the local puncture conjugacy classes are fixed by the admissible cusp
labels. Therefore \Cref{thm:d4-monodromy-rigidity} determines one monodromy representation
\[
\rho_F:\pi_1(X_f^\circ)\to SU(3)_F
\]
uniquely up to global $SU(3)_F$ conjugation.

Construct the flat Hermitian bundle explicitly as
\[
L_F:=(\widetilde X_f^\circ\times\mathbb C^3)/\pi_1(X_f^\circ),
\]
where $\pi_1(X_f^\circ)$ acts on $\mathbb C^3$ through $\rho_F$. Because the image of $\rho_F$ lies
in $SU(3)_F$, the standard Hermitian form on $\mathbb C^3$ descends to a Hermitian metric on
$L_F$, and the trivial connection on $\widetilde X_f^\circ\times\mathbb C^3$ descends to a flat
Hermitian connection realizing the same monodromy representation.

Conversely, any admissible flat Hermitian family connection with the same puncture data has the same
monodromy representation up to global $SU(3)_F$ conjugation. The unitary Riemann--Hilbert
correspondence for flat Hermitian bundles then identifies it with the descended bundle above.
Hence there is exactly one admissible flat Hermitian connection class realizing the rigid puncture
data. We denote this distinguished class by
\[
[\nabla_F^\star].
\]
\end{proof}

\begin{theorem}[Canonical transport kernel and single-winding closure]
\label{thm:exact-transport-closure}
\label{thm:canonical-transport-kernel}
The admissible cusp set is
\[
\left\{1,\frac23,\frac13\right\},
\]
obtained from the singlet-projected hypercharge spectrum. On the base hexagon the residue packets
are
\[
R_u=(0,3,1),
\qquad
R_d=(2,5,1),
\qquad
R_e=(3,5,2).
\]
Let $\mathcal G_{D_4}\subset X_f^\circ$ denote the unique $D_4$-invariant geodesic spine of the
four-punctured family sphere, and let
\[
\Gamma_{ij}^{\min}
\]
be its canonical admissible edges. Then the Yukawa matrices are fixed intrinsically by the operator
formula
\[
(Y_f)_{ij}
:=
\left\langle e_i,\mathrm{Hol}_{\mathcal G_{D_4}}(\nabla_F^\star)\,K_f^{\mathrm{str}}e_j\right\rangle,
\]
and on diagonal reduction this becomes the hard holonomy form
\[
(Y_f)_{ij}
=
\lambdaY^{L^L_{f,i}+L^R_{f,j}}
\left\langle e_i,\mathrm{Hol}_{\Gamma_{ij}^{\min}}(\nabla_F^\star)e_j\right\rangle,
\]
so the bounded coefficients $\Lambda_{f,j}$ are exact matrix entries of the structural transport
kernel rather than independent $O(1)$ amplitudes. The resonance pole of the transport kernel is the
unique admissible root
\[
\delta_{\mathrm{ph}}\in\left(\frac13,\frac23\right)
\]
of \Cref{thm:algebraic-transport-pole}. The exact sum rule is
\[
\sum_{f,j}L_{f,j}=\Oadm\gamma=40.
\]
The base residue sum is $22$, so the missing $18=3\cdot 6$ forces exactly one full winding in
each first-generation direction. Therefore uniquely
\[
L_u^{\mathrm{diag}}=(7,3,0),\quad
L_d^{\mathrm{diag}}=(7,5,2),\quad
L_e^{\mathrm{diag}}=(8,5,3)
\]
on the closed branch.
\end{theorem}

\begin{proof}
By \Cref{thm:rigid-family-local-system}, the flat Hermitian family-connection class
\[
[\nabla_F^\star]
\]
is uniquely fixed on the rigid four-punctured family sphere. By
\Cref{def:canonical-cycle-holonomy,thm:d4-monodromy-rigidity}, its holonomy on the unique
$D_4$-invariant geodesic spine $\mathcal G_{D_4}$ is therefore intrinsic. At the same time,
\Cref{def:structural-transport-kernel,thm:algebraic-transport-pole} fixes the structural positive
transport kernel and its unique admissible resonance pole
\[
\delta_{\mathrm{ph}}\in\left(\frac13,\frac23\right).
\]
Consequently the operator product
\[
\mathrm{Hol}_{\mathcal G_{D_4}}(\nabla_F^\star)\,K_f^{\mathrm{str}}
\]
is uniquely determined on the canonical family basis, and its matrix entries are theorem-level
transport data. This proves the operator formula for $Y_f$ and shows that the former coefficients
$\Lambda_{f,j}$ are not independent amplitudes but exact entries of the same closed transport
operator.

On the diagonal branch, the canonical admissible edges $\Gamma_{ij}^{\min}$ recover the word-length
reduction
\[
y_{f,j}^{\mathrm{UV}}=\lambdaY^{L_{f,j}}\Lambda_{f,j}.
\]
The admissible cusp set is the singlet-projected hypercharge packet
\[
\left\{1,\frac23,\frac13\right\},
\]
so the diagonal residue packets are exactly
\[
R_u=\{0,1,3\},
\qquad
R_d=\{1,2,5\},
\qquad
R_e=\{2,3,5\}.
\]
Write every diagonal length as
\[
L_{f,j}=r_{f,j}+6n_{f,j},
\qquad
r_{f,j}\in R_f,
\qquad
n_{f,j}\in\mathbb N_0.
\]
Summing over all sectors gives
\[
\sum_{f,j}L_{f,j}
=
\sum_{f,j}r_{f,j}+6\sum_{f,j}n_{f,j}.
\]
The residue sum is
\[
(0+1+3)+(1+2+5)+(2+3+5)=22,
\]
whereas \Cref{eq:41-chain-main} and the closed transport sum rule give
\[
\sum_{f,j}L_{f,j}=\Oadm\gamma=40.
\]
Therefore
\[
6\sum_{f,j}n_{f,j}=40-22=18,
\qquad
\sum_{f,j}n_{f,j}=3.
\]
So exactly three full windings occur in total.

The universal first-generation winding lock and the monotone hierarchy in each sector place every
extra winding on the lightest family direction. Hence each sector contributes exactly one full
winding in its first-generation slot. Reattaching the unique extra $6$ to the residue packets gives
\[
L_u^{\mathrm{diag}}=(7,3,0),\qquad
L_d^{\mathrm{diag}}=(7,5,2),\qquad
L_e^{\mathrm{diag}}=(8,5,3).
\]
This proves the closed-branch single-winding pattern.
\end{proof}

\begin{corollary}[Single-winding uniqueness on the transport branch]
\label{cor:single-winding-transport}
Writing the residue packets of \Cref{thm:exact-transport-closure} as multisets
\[
R_u=\{0,1,3\},
\qquad
R_d=\{1,2,5\},
\qquad
R_e=\{2,3,5\},
\]
the closed-branch diagonal length multisets are
\[
\{L_{u,j}^{\mathrm{diag}}\}=\{0,3,1+6\},
\qquad
\{L_{d,j}^{\mathrm{diag}}\}=\{2,5,1+6\},
\qquad
\{L_{e,j}^{\mathrm{diag}}\}=\{3,5,2+6\}.
\]
Hence each sector carries the same winding excess
\[
\sum_j L_{f,j}^{\mathrm{diag}}-\sum_{r\in R_f}r=6,
\]
and therefore the total excess is $18$. For the charged-lepton packet one also has
\[
\frac{1}{3}\sum_{j=1}^3 L_{e,j}^{\mathrm{diag}}=\frac{16}{3},
\qquad
\operatorname{Var}(L_e^{\mathrm{diag}})
=
\frac{1}{3}\sum_{j=1}^3\left(L_{e,j}^{\mathrm{diag}}-\frac{16}{3}\right)^2
=
\frac{38}{9}.
\]
\end{corollary}

\begin{proof}
\Cref{thm:exact-transport-closure} states that each sector is obtained from its residue packet by
adding exactly one full winding of length $6$ to one entry; reordering gives the displayed
multisets. The winding-excess identities are immediate. For $L_e^{\mathrm{diag}}=(8,5,3)$, direct
calculation gives the mean $16/3$ and variance $38/9$.
\end{proof}

Here the closed branch uses the exact root of \Cref{thm:transport-pole-algebraic}, so
\[
\epsilon_f
=
\frac23-\delta_{\mathrm{ph}}
\approx \num{0.073389386533443}.
\]
The later empirical readout ledgers now keep this same exact pole and only apply the declared threshold
map to downstream observables. The regularized transport operator used for the finite
hexagon analysis is
\begin{equation}
\label{eq:yukawa-kernel}
\begin{aligned}
D_y&=y\,\mathbf 1-\delta_{\mathrm{ph}}U_6,\\
\mathcal Y_y^{(\varepsilon)}&=(D_y^\dagger D_y+\varepsilon^2)^{-1},
\end{aligned}
\end{equation}
and the singular limit $\varepsilon\to0$ is taken only after admissibility has been restricted to the gap-stable sector defined below.

\begin{proposition}[Finite hexagonal resolvent]
\label{prop:finite-hexagonal-resolvent}
Whenever
\[
y^6\neq \delta_{\mathrm{ph}}^6,
\]
the transport denominator has the exact inverse
\begin{equation}
\label{eq:finite-hexagonal-resolvent}
D_y^{-1}
=
\frac{1}{y^6-\delta_{\mathrm{ph}}^6}
\sum_{m=0}^{5}y^{5-m}\delta_{\mathrm{ph}}^{\,m}U_6^m.
\end{equation}
Thus the $C_6$ resolvent is finite and algebraically closed rather than an infinite formal Green-series expansion.
\end{proposition}

\begin{proof}
Using $U_6^6=\mathbf 1$,
\[
\bigl(y\,\mathbf 1-\delta_{\mathrm{ph}}U_6\bigr)
\sum_{m=0}^{5}y^{5-m}\delta_{\mathrm{ph}}^{\,m}U_6^m
=
y^6\mathbf 1-\delta_{\mathrm{ph}}^6U_6^6
=
\bigl(y^6-\delta_{\mathrm{ph}}^6\bigr)\mathbf 1.
\]
Dividing by $y^6-\delta_{\mathrm{ph}}^6$ gives \Cref{eq:finite-hexagonal-resolvent}.
\end{proof}

Equivalently, one may read \Cref{eq:yukawa-kernel} in the Fourier basis of the six-site ring: $U_6$ is the diagonal form of the step operator, $D_y$ is the shifted walk denominator, and $\mathcal Y_y^{(\varepsilon)}$ is the corresponding regularized Green function. In this reading the admissible path lengths $L_{f,j}$ are graph-distance data on $C_6$, while the transport matrix elements $\Lambda_{f,j}$ encode bounded local amplitudes rather than unrestricted fit parameters.

Because $U_6$ has eigenphases $e^{in\pi/3}$, the regularized transport spectrum is
\begin{equation}
\label{eq:transport-spectrum}
\lambda_{y,n}^{-1}
=
y^2+\delta_{\mathrm{ph}}^2-2y\delta_{\mathrm{ph}}\cos\!\left(\frac{n\pi}{3}\right)+\varepsilon^2,
\qquad
n=0,\dots,5,
\end{equation}
for the eigenvalues $\lambda_{y,n}$ of $\mathcal Y_y^{(\varepsilon)}$. In particular,
\begin{lemma}[Finite admissible hexagon gap]
\label{lem:finite-admissible-hexagon-gap}
On the admissible cusp set
\[
y\in\left\{1,\frac13,\frac23\right\},
\]
the smallest singular value of the unregulated hexagon denominator is attained at $y=\frac23$ and $n=0$. Equivalently,
\[
\epsilon_f
:=
\min_{\substack{y\in\{1,\frac13,\frac23\}\\ n=0,\dots,5}}
\left|y-\delta_{\mathrm{ph}}e^{in\pi/3}\right|
=
\frac23-\delta_{\mathrm{ph}}.
\]
Hence the finite admissible hexagon has a strictly positive spectral gap $\epsilon_f$.
\end{lemma}

\begin{proof}
For $\varepsilon=0$, the eigenvalues in \Cref{eq:transport-spectrum} reduce to
$|y-\delta_{\mathrm{ph}}e^{in\pi/3}|^{-2}$. By \Cref{thm:transport-pole-algebraic},
$\delta_{\mathrm{ph}}$ is the unique admissible root in $(1/3,2/3)$. Hence the positive-real cusp
$y=2/3$ with $n=0$ is the nearest admissible point, while every other cusp or nonzero phase has
strictly larger distance. Therefore
\[
\epsilon_f=\frac23-\delta_{\mathrm{ph}},
\]
which is positive because $\delta_{\mathrm{ph}}<2/3$.
\end{proof}

\begin{definition}[Gap-stable admissibility]
\label{def:gap-stable-admissibility}
Let
\[
D_0:=D_{2/3}.
\]
We say that the admissible transport sector is gap-stable if the physical transport operator has the form
\[
D_f=D_0+V_{\mathrm{adm}},
\qquad
[P_{\mathrm{str}},V_{\mathrm{adm}}]=0,
\qquad
\|V_{\mathrm{adm}}\|<\frac{\epsilon_f}{2}.
\]
\end{definition}

\begin{proposition}[Structural gap lift]
\label{prop:gap-lift-full-operator}
If admissibility is gap-stable in the sense of \Cref{def:gap-stable-admissibility}, then for every
\[
\psi\in\mathrm{Ran}(P_{\mathrm{str}})
\]
one has
\[
\|D_f\psi\|
\ge
\frac{\epsilon_f}{2}\|\psi\|.
\]
Consequently
\begin{equation}
\label{eq:gap-lift-full-operator}
P_{\mathrm{str}} D_f^\dagger D_f P_{\mathrm{str}}
\ge
\frac{\epsilon_f^2}{4}\,P_{\mathrm{str}}
>
0.
\end{equation}
In particular, the final admissible restriction to $\mathrm{Ran}(\Padm)$ is strictly positive.
\end{proposition}

\begin{proof}
Since $\psi\in\mathrm{Ran}(P_{\mathrm{str}})$ and $[P_{\mathrm{str}},V_{\mathrm{adm}}]=0$,
\[
\|D_f\psi\|
\ge
\|D_0\psi\|-\|V_{\mathrm{adm}}\psi\|
\ge
\bigl(\epsilon_f-\|V_{\mathrm{adm}}\|\bigr)\|\psi\|
\ge
\frac{\epsilon_f}{2}\|\psi\|.
\]
Squaring gives \Cref{eq:gap-lift-full-operator}.
\end{proof}

\begin{equation}
\label{eq:mobius-quotient}
\frac{\lambda_{y,0}}{\lambda_{y,3}}
=
\frac{(y+\delta_{\mathrm{ph}})^2+\varepsilon^2}{(y-\delta_{\mathrm{ph}})^2+\varepsilon^2}.
\end{equation}
\[
\lim_{\varepsilon\to 0}\frac{\lambda_{y,0}}{\lambda_{y,3}}
=
\left(\frac{y+\delta_{\mathrm{ph}}}{y-\delta_{\mathrm{ph}}}\right)^2
\]
\begin{remark}[Near-pole hierarchy mechanism]
On the closed transport branch, the exact pole is the root of
\Cref{prop:closed-form-transport-pole}; the later comparison ledgers retain that same pole and
regularize only the finite Green kernel at $\varepsilon>0$.
Hence the Möbius quotient contains a built-in near-pole whenever $y$ approaches $\delta_{\mathrm{ph}}$, making large hierarchy ratios structurally plausible. In the closed transport grammar this is the hierarchy mechanism, while the bounded holonomy factors keep it from degenerating into an arbitrary tuning surface.
\end{remark}
